预处理在 Croqtile 解析之前进行。普通 C++ 和 \co@__co__@ 区域可以共用宏定义，便于统一维度与配置。

\begin{table}[htbp]\centering\small
\begin{tabularx}{\linewidth}{@{}lX@{}}
\toprule
指令 & 作用 \\
\midrule
\texttt{\#define NAME value} & 对象式宏：文本替换 \\
\texttt{\#if} / \texttt{\#elif} / \texttt{\#else} / \texttt{\#endif} & 条件包含 \\
\texttt{\#ifdef} / \texttt{\#ifndef} & 判断宏是否已定义的简写 \\
\texttt{\#error message} & 带消息强制编译时失败 \\
\bottomrule
\end{tabularx}
\end{table}

预处理器不支持函数式宏。需要带参数的计算时，可以在普通 C++ 中使用辅助函数。
